\chapter{Conclusion and Future Work}
\label{chap:conclusion}

This chapter summarizes the achievements of the L3S-Offshore-2 Frontend project and outlines directions for future development.

\section{Summary}
\label{sec:summary}

This documentation has described the development of a web-based frontend for Petri net visualization and simulation, created as part of the L3S-Offshore-2 research project.

\subsection{Achievements}

The following goals were accomplished:

\begin{enumerate}
    \item \textbf{Functional Visualization:} The application successfully renders Petri nets from PNML files with interactive editing capabilities.
    
    \item \textbf{Simulation Integration:} Full integration with the Flask backend enables multi-run simulations with firing sequence animation.
    
    \item \textbf{Offshore Parameter Interface:} A dedicated tab allows researchers to configure domain-specific parameters for scheduling models.
    
    \item \textbf{Production Deployment:} The application is containerized and deployed on L3S infrastructure, accessible via web browser.
    
    \item \textbf{Code Quality:} Extensive JSDoc documentation and English comments ensure maintainability for future developers.
\end{enumerate}

\subsection{Technical Highlights}

The most technically challenging aspects that were solved:

\begin{itemize}
    \item \textbf{Viewport-centered zoom:} Implementing intuitive zoom behavior with ResizeObserver integration
    \item \textbf{Tab-state management:} Ensuring UI consistency when switching between different workflows
    \item \textbf{Token reset logic:} Preserving initial markings across simulation runs
    \item \textbf{Reactive forms integration:} Solving performance issues with Angular's form system
\end{itemize}

\section{Limitations}
\label{sec:limitations}

The current implementation has the following limitations:

\begin{itemize}
    \item \textbf{No data persistence:} All data is stored in-memory and lost on page reload
    \item \textbf{Limited to P/T nets:} Higher-level Petri net types (colored, hierarchical) are not supported
    \item \textbf{Single-user:} No collaborative editing or user accounts
    \item \textbf{Visual overlaps:} Large nets may have overlapping elements with the current layout
\end{itemize}

\section{Future Work}
\label{sec:future-work}

The following enhancements could be pursued in future development:

\subsection{Short-Term (High Priority)}

\begin{itemize}
    \item \textbf{Interactive Miner Integration:} Connect to pm4py's interactive miner for process discovery
    \item \textbf{Analysis Endpoints:} Display conformance checking results from the backend
    \item \textbf{Export Functions:} Enable export of simulation results as CSV/JSON
\end{itemize}

\subsection{Medium-Term (Nice-to-Have)}

\begin{itemize}
    \item \textbf{Undo/Redo:} Implement history management for the Build tab
    \item \textbf{Data Persistence:} Add optional browser-based storage (IndexedDB)
    \item \textbf{Visual Improvements:} Reduce element overlaps with improved layout algorithms
\end{itemize}

\subsection{Long-Term (Research Extensions)}

\begin{itemize}
    \item \textbf{Colored Petri Nets:} Extend the data model and visualization for CPNs
    \item \textbf{Real-Time Monitoring:} WebSocket-based live updates during long-running simulations
    \item \textbf{Multi-User Collaboration:} Shared editing with conflict resolution
\end{itemize}

\section{Acknowledgements}
\label{sec:acknowledgements}

This project was developed at the L3S Research Center, Leibniz Universität Hannover, under the supervision of Peng Qian. The author thanks the original developers of the i-love-petrinets project for providing the initial codebase.